% Generisano: uređuj beleske/sNNN.tex i naslove u manifestu.
\documentclass[11pt,a4paper]{article}
% Zajednički preambulum za dokument za čitanje; nije Beamer tema.
\usepackage[a4paper,margin=20mm,headheight=15pt]{geometry}
\usepackage{lmodern,fontspec}
\usepackage[serbian]{babel}
\setmainfont{Latin Modern Sans}
\setsansfont{Latin Modern Sans}
\setmonofont{Latin Modern Mono}
\usepackage{amsmath,amssymb,mathtools,graphicx,xcolor}
\usepackage{array,booktabs,tabularx,multirow,makecell,listings,fancyhdr}
\pagestyle{fancy}\fancyhf{}
\newcommand{\BeleskeTekuciID}{NEOZNACEN}
\fancyhead[L]{\small Beleške uz slajd \texttt{\expandafter\detokenize\expandafter{\BeleskeTekuciID}}}
\fancyfoot[R]{\small\thepage}
\setlength{\parindent}{0pt}\setlength{\parskip}{6pt}
\newwrite\BeleskeMapa
\immediate\openout\BeleskeMapa=\jobname.notes-map.tsv
\AddToHook{shipout/before}{%
  \immediate\write\BeleskeMapa{\BeleskeTekuciID\space\thepage}}
\AddToHook{enddocument/afterlastpage}{\immediate\closeout\BeleskeMapa}
\newcommand{\BeleskeID}[1]{%
  \def\BeleskeZadatiID{#1}%
  \ifx\BeleskeTekuciID\BeleskeZadatiID\else
    \PackageError{beleske}{Pogresan ID beleske: #1}{Proveri mapiranje slajda.}%
  \fi}
\newcommand{\BeleskaSlajda}[4]{%
  \clearpage\gdef\BeleskeTekuciID{#1}%
  {\Large\bfseries Slajd \textnormal{\texttt{\detokenize{#1}}}: #3\par}\medskip
  \includegraphics[page=#2,width=\linewidth]{\BeleskePrezentacija}\par\medskip
  \input{#4}}

\newcommand{\BeleskePrezentacija}{build/09_2_mos_logicka_kola.pdf}
\begin{document}
\BeleskaSlajda{09_2-s001}{1}{Logička kola sa MOS tranzistorima — 2. deo}{beleske/s001.tex}
\BeleskaSlajda{09_2-s002}{2}{NMOS invertor sa ugrađenim kanalom}{beleske/s002.tex}
\BeleskaSlajda{09_2-s003}{3}{Početak provođenja pogonskog tranzistora}{beleske/s003.tex}
\BeleskaSlajda{09_2-s004}{4}{Donji ulazni prag: uslov nagiba −1}{beleske/s004.tex}
\BeleskaSlajda{09_2-s005}{5}{Oba NMOS tranzistora u zasićenju}{beleske/s005.tex}
\BeleskaSlajda{09_2-s006}{6}{Veliko pojačanje i tačka prebacivanja}{beleske/s006.tex}
\BeleskaSlajda{09_2-s007}{7}{Gornji ulazni prag NMOS invertora}{beleske/s007.tex}
\BeleskaSlajda{09_2-s008}{8}{Napon logičke nule}{beleske/s008.tex}
\BeleskaSlajda{09_2-s009}{9}{Održivost naponskih nivoa}{beleske/s009.tex}
\BeleskaSlajda{09_2-s010}{10}{Karakteristika prenosa: ugrađeni kanal}{beleske/s010.tex}
\BeleskaSlajda{09_2-s011}{11}{Invertor u pseudo NMOS tehnologiji}{beleske/s011.tex}
\BeleskaSlajda{09_2-s012}{12}{PMOS: negativna konvencija, triodna oblast}{beleske/s012.tex}
\BeleskaSlajda{09_2-s013}{13}{PMOS: negativna konvencija, zasićenje}{beleske/s013.tex}
\BeleskaSlajda{09_2-s014}{14}{PMOS: pozitivna konvencija, triodna oblast}{beleske/s014.tex}
\BeleskaSlajda{09_2-s015}{15}{PMOS: pozitivna konvencija, zasićenje}{beleske/s015.tex}
\BeleskaSlajda{09_2-s016}{16}{Pseudo NMOS: napon logičke jedinice}{beleske/s016.tex}
\BeleskaSlajda{09_2-s017}{17}{Pseudo NMOS: određivanje donjeg praga}{beleske/s017.tex}
\BeleskaSlajda{09_2-s018}{18}{Pseudo NMOS: pragovi i skaliranje širina}{beleske/s018.tex}
\BeleskaSlajda{09_2-s019}{19}{Pseudo NMOS: oba tranzistora u zasićenju}{beleske/s019.tex}
\BeleskaSlajda{09_2-s020}{20}{Pseudo NMOS: gornji ulazni prag}{beleske/s020.tex}
\BeleskaSlajda{09_2-s021}{21}{Pseudo NMOS: napon logičke nule}{beleske/s021.tex}
\BeleskaSlajda{09_2-s022}{22}{Pseudo NMOS: naponski nivoi i odnos beta}{beleske/s022.tex}
\BeleskaSlajda{09_2-s023}{23}{Invertor u CMOS tehnologiji}{beleske/s023.tex}
\BeleskaSlajda{09_2-s024}{24}{CMOS: izlazna logička jedinica}{beleske/s024.tex}
\BeleskaSlajda{09_2-s025}{25}{CMOS: N u zasićenju, P u omskoj oblasti}{beleske/s025.tex}
\BeleskaSlajda{09_2-s026}{26}{CMOS: ravnoteža struja u drugoj oblasti}{beleske/s026.tex}
\BeleskaSlajda{09_2-s027}{27}{CMOS: aproksimacija donjeg ulaznog praga}{beleske/s027.tex}
\BeleskaSlajda{09_2-s028}{28}{CMOS: donji prag simetričnog invertora}{beleske/s028.tex}
\BeleskaSlajda{09_2-s029}{29}{CMOS: oba tranzistora u zasićenju}{beleske/s029.tex}
\BeleskaSlajda{09_2-s030}{30}{CMOS: veliko pojačanje u trećoj oblasti}{beleske/s030.tex}
\BeleskaSlajda{09_2-s031}{31}{CMOS: N u omskoj oblasti, P u zasićenju}{beleske/s031.tex}
\BeleskaSlajda{09_2-s032}{32}{CMOS: ravnoteža u četvrtoj oblasti}{beleske/s032.tex}
\BeleskaSlajda{09_2-s033}{33}{CMOS: izvođenje praga u četvrtoj oblasti}{beleske/s033.tex}
\BeleskaSlajda{09_2-s034}{34}{CMOS: gornji prag simetričnog invertora}{beleske/s034.tex}
\BeleskaSlajda{09_2-s035}{35}{CMOS: izlazna logička nula}{beleske/s035.tex}
\BeleskaSlajda{09_2-s036}{36}{CMOS: jednake margine šuma}{beleske/s036.tex}
\BeleskaSlajda{09_2-s037}{37}{CMOS: pet oblasti karakteristike}{beleske/s037.tex}
\BeleskaSlajda{09_2-s038}{38}{Podešavanje tačke prebacivanja CMOS invertora}{beleske/s038.tex}
\BeleskaSlajda{09_2-s039}{39}{Tačka prebacivanja: dugačak kanal}{beleske/s039.tex}
\BeleskaSlajda{09_2-s040}{40}{Dimenzionisanje za željenu tačku prebacivanja}{beleske/s040.tex}
\BeleskaSlajda{09_2-s041}{41}{Tačka prebacivanja: kratak kanal}{beleske/s041.tex}
\BeleskaSlajda{09_2-s042}{42}{Odnos širina: model kratkog kanala}{beleske/s042.tex}
\BeleskaSlajda{09_2-s043}{43}{Kratak kanal: prag na polovini napajanja}{beleske/s043.tex}
\BeleskaSlajda{09_2-s044}{44}{Tačka prebacivanja: izraženo zasićenje brzine}{beleske/s044.tex}
\BeleskaSlajda{09_2-s045}{45}{Odnos širina u linearnom modelu}{beleske/s045.tex}
\BeleskaSlajda{09_2-s046}{46}{Normalizovane širine tranzistora}{beleske/s046.tex}
\BeleskaSlajda{09_2-s047}{47}{Jedinični CMOS invertor i tehnološke dimenzije}{beleske/s047.tex}
\BeleskaSlajda{09_2-s048}{48}{Minimalna geometrija i normalizacija lambda}{beleske/s048.tex}
\BeleskaSlajda{09_2-s049}{49}{Zaštita od elektrostatičkog pražnjenja (ESD)}{beleske/s049.tex}
\BeleskaSlajda{09_2-s050}{50}{Latch-up: parazitni tranzistori}{beleske/s050.tex}
\BeleskaSlajda{09_2-s051}{51}{Latch-up: pozitivna povratna sprega}{beleske/s051.tex}
\BeleskaSlajda{09_2-s052}{52}{Latch-up: posledice i sprečavanje}{beleske/s052.tex}
\BeleskaSlajda{09_2-s053}{53}{Statička disipacija CMOS kola}{beleske/s053.tex}
\BeleskaSlajda{09_2-s054}{54}{Dva izvora dinamičke disipacije}{beleske/s054.tex}
\BeleskaSlajda{09_2-s055}{55}{Energija punjenja i pražnjenja}{beleske/s055.tex}
\BeleskaSlajda{09_2-s056}{56}{Dinamička snaga i učestanost prebacivanja}{beleske/s056.tex}
\BeleskaSlajda{09_2-s057}{57}{Napajanje, kašnjenje i disipacija}{beleske/s057.tex}
\BeleskaSlajda{09_2-s058}{58}{Struja kratkog spoja tokom ulaznih ivica}{beleske/s058.tex}
\BeleskaSlajda{09_2-s059}{59}{Energija i snaga kratkospojne struje}{beleske/s059.tex}
\BeleskaSlajda{09_2-s060}{60}{Ukupna disipacija i faktor PDP}{beleske/s060.tex}
\BeleskaSlajda{09_2-s061}{61}{Faktor EDP: energija i kašnjenje}{beleske/s061.tex}
\BeleskaSlajda{09_2-s062}{62}{Napajanje koje minimizuje EDP}{beleske/s062.tex}
\end{document}
